%% kapa: exam=1 round=19 year=2008 period=1 form=A subject=appraisal_law subject_ko=부동산관계법규 seq=37 number=77
%% source: Q-Net 감정평가사 기출문제 (cst00309) artlSeq=1001962 file=extracted/1008/1001962/2008년 제19회 감정평가사 1차시험 기출문제/1교시.hwp
%% conversion: kordoc 4.24.51 -> tools/gen1.py (auto); review=proofread-sonnet; flags=none
\begin{question}{77}
부동산등기법상 등기한 권리의 순위에 관한 설명으로 옳지 않은 것은?
\begin{choices}
  \item 동일한 부동산에 관하여 등기한 권리의 순위는 법률에 다른 규정이 없는 때에는 등기한 순서에 따른다.
  \item 등기의 순서는 등기용지상의 같은 구에서 한 등기에 대해서는 순위번호에 의하고 다른 구에서 한 등기에 대해서는 접수번호에 의한다.
  \item 부기등기 상호간의 순위는 주등기의 순위에 의한다.
  \item 가등기를 한 경우에 본등기의 순위는 가등기의 순위에 의한다.
  \item 부동산강제경매절차에서 선순위의 담보권이나 가압류가 없는 소유권이전등기청구권의 순위보전을 위한 가등기가 담보가등기인지 순위보전의 가등기인지 밝혀질 때까지 경매절차를 중지해야 하는 것은 아니다.
\end{choices}
\end{question}
\answer{3}
